% This file was converted to LaTeX by Writer2LaTeX ver. 1.9.9
% see http://writer2latex.sourceforge.net for more info
\documentclass{article}
\usepackage{calc,amsmath,amssymb,amsfonts}
\usepackage[LGR,T1]{fontenc}
\usepackage[greek,italian]{babel}
\usepackage[style=numeric,backend=biber]{biblatex}
\author{HP}
\date{2026-08-19}
\begin{document}
Vincenzo Camporeale è informatico, programmatore e appassionato di retrocomputing.

Il suo incontro con l'informatica risale ai primi anni Ottanta, quando, ancora bambino, davanti a una partita a Pong
giocata a casa di un amico si accorse che a incuriosirlo non era soltanto ciò che accadeva sullo schermo, ma
soprattutto ciò che non poteva vedere.

Qualche anno dopo incontrò il Commodore 64 e il BASIC durante un corso di programmazione tenuto dal professor Rossi
Brunori. Poco dopo, i suoi genitori gli regalarono il suo primo computer: proprio un Commodore 64, destinato a
contribuire in modo decisivo a orientare i suoi studi e, molti anni più tardi, la sua vita professionale.

Come molti ragazzi di quella generazione imparò anche copiando e modificando i listati pubblicati sulle riviste,
acquistava regolarmente Commodore Computer Club $\text{\textgreek{—}}$ che arrivò a pubblicare una sua lettera dedicata
alla possibilità di utilizzare la user port per comandare un videoregistratore $\text{\textgreek{—}}$ e trascorreva
davanti alla tastiera almeno tanto tempo a cercare di capire la macchina quanto a giocare.

Diplomato in Informatica e successivamente laureato in Informatica presso l'Università degli Studi di Bari, lavora da
circa trent'anni nel settore dello sviluppo software e dei sistemi informativi. Nel corso della sua attività
professionale ha attraversato tecnologie e paradigmi molto diversi, dai sistemi legacy e Progress 4GL/ABL allo sviluppo
web, fino ad ASP.NET Core e alle moderne architetture applicative, occupandosi anche di database, migrazione dei dati,
integrazione tra sistemi e interoperabilità con piattaforme esterne.

Alla programmazione affianca da sempre una forte curiosità per ciò che accade dietro le astrazioni del software: capire
i meccanismi, sperimentare, mettere alla prova le proprie convinzioni e ricostruire il percorso che conduce da un
problema alla sua soluzione. Una curiosità che coltiva anche attraverso il retrocomputing, l'elettronica e il Commodore
64, tornando spesso alle radici dell'informatica per comprenderne meglio il presente.

Circa quarant'anni dopo quei primi esperimenti, avvicinandosi nuovamente al mondo del retrocomputing, ha ritrovato
proprio il suo primo computer, quel Commodore 64 che i suoi genitori gli avevano regalato affrontando un sacrificio
economico importante. La macchina originale non era più funzionante, ma invece di sostituirla ha scelto di conservarne
il case e la tastiera e di riportarla progressivamente in vita attraverso un progetto basato su Raspberry Pi e VICE,
evitando di alterarne irreversibilmente le parti originali.

Il ritorno al C64 ha riaperto anche un'altra storia rimasta incompiuta: il desiderio, nato da ragazzo insieme agli amici
di allora, di creare un videogioco tutto loro. Questo libro nasce dal tentativo di mantenere finalmente quella promessa
e, nello stesso tempo, di condividere con il lettore il modo di ragionare di una generazione che spesso ha imparato
l'informatica proprio così: davanti a un cursore lampeggiante, con poche risorse e molta sperimentazione.

Vive in Puglia con la sua famiglia e continua a programmare, studiare e sperimentare con la stessa curiosità con cui, da
bambino, cercava di capire che cosa ci fosse dietro lo schermo di un Commodore 64.
\end{document}
